% An independent structural certificate. Previous bounded-search UNKNOWN
% outputs are retained byte for byte and are not used in this proof.
\begin{lemma}
The listed pentagon component with boundary $00001$ has no Hamilton path.
\end{lemma}
\begin{proof}
All four prefix colours of a state $x$ lie in $\{1,2,3\}$, and
$x_4\in\{0,2,3\}$. Orient the colour cycle $1\to2\to3\to1$,
put $o(a,b)=1$ for its forward edges and $o(a,b)=-1$ for its reverse
edges, and set $\sigma(x)=o(x_0,x_1)o(x_2,x_3)$.
Let $\chi(x)=\sigma(x)$ if $x_4=0$, and $\chi(x)=-\sigma(x)$ otherwise.
A legal move in the first four coordinates changes the sign of $\sigma$.
A move from $0$ to $2$ or $3$ in the last coordinate changes the sign
of $\chi$. The only remaining possible moves have $x_4$ changing between
$2$ and $3$; they require $x_0=x_3=1$ and hence are precisely
$12312\,12313$ and $13212\,13213$. Both have $\chi=-1$.

The counts for $(x_4,\sigma)$ are $12,12$ at $x_4=0$ and $6,4$ at
each of $x_4=2,3$. Thus $S=\{x:\chi(x)=1\}$ has twenty vertices.
The induced graph on the other twenty-four vertices has exactly the
two disjoint edges just displayed, so $c(B-S)=22>|S|+1=21$.
Removing twenty vertices from a Hamilton path leaves at most twenty-one
path components. Therefore $B$ has no Hamilton path.
\end{proof}
